\section{Experimental setup}
\label{sec:setup}

\subsection{Field, dynamics and reference}

\ph{one paragraph} GRAIL product and degree \citep{Konopliv2014GRAIL,
GrailGL1800FSHADR,Goossens2020}; synthesis kernel and normalization
\citep{holmes2002unified,pines1973uniform}; perturbation model and the
conventions it follows \citep{montenbruck2000satellite,vallado2013fundamentals};
third bodies and ephemeris \citep{park2021de440,acton2018spice}; integrator and
tolerances \citep{hairer1993ode,virtanen2020scipy}; lunar orientation
\citep{Mazarico2018}. Adopted \emph{reference} degree per orbit and the sense in
which every result below is model-relative: what is measured is what the adopted
expansion requires of its own truncations, not the distance to the true lunar
field.

\subsection{Populations}

\begin{table}[!htbp]
  \centering\small
  \caption{Orbit populations. Designs A and B are independent scrambles of the
  same scrambled-Sobol coverage construction \citep{joe2008sobol}; C was drawn
  and frozen after A and B were known to agree. That the design points are
  dependent by construction is why no significance level is attached to any
  tally. The wide-elliptic population lies outside the sampled factor box
  and brackets the elliptical subsatellite orbits flown on Kaguya
  \citep{kato2010kaguya}.}
  \label{tab:populations}
  \begin{tabular}{@{}llll@{}}
    \toprule
    Population & Orbits & Geometry & Role \\
    \midrule
    Design A & \ph{n} & \ph{box} & primary \\
    Design B & \ph{n} & \ph{box} & independent replication \\
    Design C & \ph{n} & \ph{box} & third design \\
    Wide-elliptic & \ph{n} & perilune \ph{km}, apolune \ph{km} & regime bound \\
    Strata (5) & \ph{n} each & inclination / argument of perilune restricted & geometry bound \\
    Low perilune & \ph{n} & \ph{km} & linearization bound \\
    \bottomrule
  \end{tabular}
\end{table}

\subsection{Budgets}
\label{sec:setup-budgets}

Four cost quantities appear and are never interchanged.

\begin{table}[!htbp]
  \centering\small
  \caption{Cost definitions. $B_2$ is the primary budget throughout; $B_1$ is
  reported only for comparability with the previous paper.}
  \label{tab:budget-definitions}
  \begin{tabular}{@{}llp{0.44\textwidth}@{}}
    \toprule
    Symbol & Quantity & Note \\
    \midrule
    $B_1$ & nominal work rate, the \emph{time average} $\langle N^2\rangle_t
            = \sum_i N_{g(i)}^2\Delta t_i/T$
          & available without propagating; does not survive contact with the
            integrator. On the uniform reference output grid this is the plain
            mean, which is how the previous paper reported it \\
    $B_2$ & realized total quadratic work $\sum_{\mathrm{RHS}}N^2$
          & \textbf{primary}; requires propagation \\
    $B_3$ & measured serial gravity-kernel time
          & architecture-dependent; idle machine, three repeats \\
    $B_+$ & controller overhead, in the same units as $B_2$
          & \textbf{included in every budget reported for \Cplan, \Clite\ and
            \Cfb}; conversion below \\
    \bottomrule
  \end{tabular}
\end{table}

\paragraph{What $B_+$ is measured in.} $B_2$ counts quadratic work,
$\sum_{\mathrm{RHS}}N^2$, and $B_+$ is added to it, so the two must carry the
same units. Two of the overhead terms already do: the pilot arc and the online
probe are gravity evaluations and enter directly as $\sum N^2$. The rest ---
the two-body predictor, the planning sweeps, the online $\argmin$ and its
bookkeeping --- are processor time and are converted at the measured rate of
the same run,

\begin{equation}
  B_{+,\mathrm{eq}} \;=\; t_+ \Bigl/ \frac{t_{\mathrm{kernel}}}{B_2} ,
  \label{eq:overhead-units}
\end{equation}

with $t_{\mathrm{kernel}}$ the arc's measured serial gravity-kernel time. The
ratio $t_{\mathrm{kernel}}/B_2$ is a seconds-per-unit-work coefficient, so
Eq.~\eqref{eq:overhead-units} expresses non-gravity time in $N^2$ units and
Eq.~\eqref{eq:matching} adds like to like. The coefficient is
architecture-dependent and is measured under the same conditions as $B_3$, on
an idle machine with three repeats; every table quoting $B_+$ names the machine
it was measured on.

\paragraph{Normalization.} $\beta$ is defined at the $B_1$ level only, as the
allocation \emph{target} $\beta=B_1/N_{\mathrm{crit}}^2$. The realized
counterpart is reported separately as
$\beta_2 = B_2/(K_{\mathrm{ref}}N_{\mathrm{crit}}^2)$, with
$K_{\mathrm{ref}}$ the comparator's right-hand-side call count on the same arc;
dividing $B_2$ by $N_{\mathrm{crit}}^2$ alone would fold the call count into
the same number and is not done.

\paragraph{Matching convention.} Comparisons are matched on a common work
\emph{ceiling}, and the controller's overhead comes out of its own allocation:

\begin{equation}
  B_2 + B_+ \;\le\; B_{\mathrm{tot}} \;=\; B_{2,\mathrm{F}} ,
  \label{eq:matching}
\end{equation}

with the inequality meant literally. Section~\ref{sec:algorithm-descent}
explains why an allocation can be better for spending less --- the objective
is not monotone in degree --- so the realized $B_2+B_+$ is reported beside
every error rather than assumed equal to the ceiling. A policy that wins under
a lower spend has won by more; one that wins only by saturating the ceiling is
reported as such.

with $B_+=0$ for every policy that is not a controller. No fictitious overhead
is added to the fixed comparator, so a controller that must plan and probe has
strictly less budget left for gravity evaluation than its comparator does ---
the conservative reading, and the one adopted throughout.

Equation~\eqref{eq:matching} is not complete until it says which side moves,
and the answer is not the one the previous campaign used. There, the comparator
was re-propagated to spend what the candidate had spent. That cannot be done
here, because the capture fraction of Section~\ref{sec:campaign} divides one
difference by another and both must be taken against \emph{the same}
$E_{\Fop}$; a comparator that slides to meet each candidate would put two
different $\Fop$ instances in the numerator and the denominator.

We therefore anchor. For each orbit and each $\beta$, $B_{\mathrm{tot}}$ is
declared in advance as the realized total quadratic work of $\Fop(\beta)$ on
that arc --- at $\beta=1$ that is the archived critical-degree run, so the
anchor already exists and is not recomputed --- and \emph{every candidate is
calibrated to it}.

That calibration costs propagations, and the cost is a consequence of the
previous paper's own result rather than an inefficiency: a schedule's realized
work is not known until it is flown, and calibrating on nominal per-call work
misses it by a median \ph{pct}. Each candidate is therefore propagated,
measured, its multiplier adjusted, and propagated again until it sits under
the ceiling without slack it does not want:
$B_2+B_+\le B_{\mathrm{tot}}$, and either
$(B_{\mathrm{tot}}-B_2-B_+)/B_{\mathrm{tot}}\le\ph{pct}$ or the schedule is
the unconstrained optimum of its own criterion. The second branch is what
keeps the calibration from forcing work onto a policy whose objective would
worsen for taking it. That takes \ph{value} propagations per candidate in
practice, and Section~\ref{sec:campaign} reports the achieved spend alongside
every verdict.

The propagated grid is $\beta\in\{0.50,0.75,1.00,1.25,1.50\}$;
$\beta=3$ is evaluated at benchmark level only, because at that budget the
previous campaign found \ph{n} of \ph{n} comparisons undecidable and the
experiment stops being an instrument.

\subsection{Grids and their convergence}
\label{sec:setup-grids}

\ph{one paragraph plus table} The campaign runs on an accumulation grid
refined on $\tau_{\mathrm{corr}}$, so the convergence parameter is not an
absolute step but the number of samples per correlation time,
$n_s=\tau_{\mathrm{corr}}(t)/\Delta t_i$, swept over
$\{0.5,1,2,4,8\}$. A uniform-grid sweep at
$\Delta t_{\mathrm{acc}}\in\{240,120,60,30,10\}$~s is run alongside it, both
because it is the comparable quantity in the previous paper and because it
isolates whether the refinement, rather than the resolution, is what matters.
Three quantities are checked in each: the objective, the recovered schedule
itself as a degree-profile distance, and verdict stability. The decision grid
is swept separately at $\Delta t_{\mathrm{dec}}\in\{60,120,300\}$~s. The
campaign runs at $n_s=\ph{value}$. \wpref{WP4}

\dnote{Section~\ref{sec:problem} predicts that the signed integral needs a
finer grid than the magnitude statistic did. If the panel contradicts that ---
if 120~s is already converged for $J$ and for the schedule --- say so plainly
and delete the prediction rather than keeping it as motivation.}

\subsection{Resolution rule and statistics}

A comparison is decided only when the two errors differ by more than their
summed reference-inclusive numerical envelope, $M_{\mathrm{res}}>1$; comparisons
below that are \emph{undecided} and are counted for neither side. No
significance level is attached to any tally: the points of a scrambled design
are dependent by construction and the resolution rule censors the sample
non-randomly, so a binomial test would have no valid null. Direction rests on
replication across independent populations. Ratios are medians of per-orbit
ratios, never ratios of medians.

\subsection{Comparators}
\label{sec:setup-comparators}

Two constant-degree comparators are carried, and they answer different
questions.

\begin{itemize}
  \item \Fop, the integer degree nearest the budget, is what a user can select
    in advance. Beating it is the claim.
  \item \Fenv, the lowest-error member of the constant family under the same
    budget, is found by measuring every candidate against the reference. It is
    a post-hoc lower envelope, not a policy; the previous paper measured it a
    median \ph{value} times better than the saturating degree. It is reported
    to bound how much of the constant family's potential \Fop\ leaves unused,
    and no hypothesis of this paper requires beating it.
\end{itemize}

\Rint, the interior member of the radial--constant interpolation family, is
carried as the strongest previously available deployable schedule.

\subsection{Pre-registration and its stopping rules}
\label{sec:setup-prereg}

Hypotheses, gates, sweeps, comparator rules, censoring and the naming rule of
Section~\ref{sec:algorithm-certificate} were written to
\ph{oa01/oa02/oa03 json} and hashed before any aggregate result was inspected.
Anything added later is labelled post hoc and reported separately; rejected
hypotheses are reported as rejected rather than reworded. \ph{outcome list}

Two of those rules bound the campaign rather than a single number, and both are
stated here because later sections invoke them.

\paragraph{Budget gate.} The first propagated campaign is at $\beta=1$ on one
design. The budget grid is extended to the rest of the populations only if the
controller also takes the resolved majority at $\beta=0.75$ \emph{and}
$\beta=1.25$ on a stratified sixteen-orbit subset. If it does not, the gain
belongs to one budget point, the population expansion is not run, and the
result is reported as such.

\paragraph{Screening escape.} The variational screen of
Section~\ref{sec:campaign} is calibrated on policies that were not built by
minimizing $J$, so its transfer to policies that were is an assumption rather
than a measurement. If the propagated sign disagrees with the screened sign,
the implementation is diagnosed; after two diagnostic rounds without
resolution the instrument is declared uncalibrated for this class, the screen's
selection is discarded, and the campaign proceeds on propagation alone with a
reduced candidate list. The cost of that path is reported rather than absorbed.
The rule exists so that a disagreement cannot become an open-ended loop.
